\section{Configurable and JD-Driven Screening}
\label{sec:lit-extensibility}

The majority of ATS platforms still view screening as a static matching problem, using the same dimensions
and weightings across roles, even when job descriptions place emphasis on different skills, tech stacks
and seniority levels. This has led to research on automated job-candidate matching, pivoting towards
a richer representation of skills, however, this has not relaxed the rigidity of the scoring system.

We begin with Gugnani and Misra~\cite{gugnani2020}, who extract both explicit and implicit skills from JDs through
large-scale document embeddings, achieving this through the training of a vector embedding model. The dataset consists
of roughly one million job postings and aims to capture the semantic relationships between a role and its
required competencies. This helps to surface skills that a recruiter may have not listed explicitly.
While this does improve the differentiation of skills, the issue still persists that the matcher relies on
a fixed set of weightings. Explicit skills, similarity, frequency, and implicit-skill penalties are combined
through hardcoded coefficients. It is evident that adaptability across occupations is bounded by the 
training corpus and also the chosen scoring formula. It is not influenced by the recruiter's job description
or priorities for a particular hire.

A different approach is taken by Lo et al.~\cite{lo2025aihiring}, who split CV screening into modular LLM stages.
The stages are as follows: resume extraction, evaluation, summarisation, and score formatting.
They integrate retrieval-augmented generation (RAG) to allow the evaluator to draw on external knowledge such
as certifications and university rankings, rather than just the CV alone.

Their evaluation study shows that the pipeline achieves a strong agreement with expert judgement, outperforming 
simpler baselines, leading to the conclusion that modular LLM workflows can approximate human screening.
The issue however, is that evaluation dimensions remain rigid, with the following fixed metrics:
skills, experience, education, background, and summary. Note that these metrics also have fixed weights,
furthering the rigidity of the scoring system. This limits the ability for the scoring system to adapt to
role-specific priorities, alternative career paths and criteria that fall outside of the standard.

Together, the studies show that both semantic and LLM-based screening have the capacity to improve accuracy and 
relevance, but still fail to address the extensibility problem. There is no answer to which metrics matter for 
a particular role, nor do they suggest weights based off the context of the job description.

We review the work of Karen Sparck Jones~\cite{sparckjones1972idf}, who proposed Term Frequency-Inverse 
Document Frequency (TF-IDF) as a way to rank the importance of terms within a document. Here, TF-IDF ranks
the importance of a term based off its appearance within a document compared to its appearance throughout a 
corpus of documents. If a term is emphasised heavily in a document relative to the corpus, it will receive a 
higher ranking due to its rarity within the corpus. In a similar fashion, commonly used terms are down-weighted
due to their abundance within the corpus. In a similar fashion to large-scale skill inference from professional 
documents~\cite{bastian2014linkedin}, corpus-level frequency has been used to separate role-specific 
requirements from generic boilerplate. This approach is adapted and applied to MERIT later during Chapter~\ref{ch:implementation}.

\section{Semantic Intelligence and Vector Embeddings}

MERIT leverages the transition from traditional lexical analysis to semantic intelligence to evaluate the technical expertise of a candidate.
Earlier generations of ATS relied on "keyword matching", however, MERIT will leverage Large Language Models (LLMs) and vector embeddings to build relationships between programming languages and technical frameworks.

\subsection*{Lexical vs. Semantic Analysis}
Traditional lexical methods operate on the principle of exact character matching. 
These methods are often efficient for identifying rare terms, such as a unique framework.
However, they are vulnerable to the "vocabulary mismatch" problem, where a candidate may not use the exact same terminology as the job description.
Failing to match a candidate due to a mismatch in terminology is a significant issue that can lead to qualified candidates being overlooked simply because they didn't use the same keywords as the job description.

Semantic models, on the other hand, represent technical skills as vectors in a high-dimensional latent space. 
By mapping technologies to coordinates, MERIT can measure the proximity of skills rather than the frequency of words. 
This allows the system to understand that these terms are conceptually related, enabling "soft" matching for candidates who possess highly transferable technical skills.

\subsection*{Vector Similarity Measures}
The similarity between a candidate's technical profile ($\mathbf{C}$) and the job requirements ($\mathbf{J}$) is calculated using measures such as Cosine Similarity:
\begin{equation}
    \text{similarity} = \cos(\theta) = \frac{\mathbf{C} \cdot \mathbf{J}}{\|\mathbf{C}\| \|\mathbf{J}\|}
\end{equation}

This mathematical foundation enables MERIT to identify related competencies and technical potential that would be missed by rigid rule-based systems. 
By utilising transformer-based architectures such as Sentence-BERT \cite{reimers2019sentence}, the system can generate dense vector representations that capture the subtle linguistic nuances of technical expertise and framework relationships.

